%
% 14-matrix.tex -- Matrixschreibweise für komplexe Zahlen
%
% (c) 2026 Prof Dr Andreas Müller
%
\subsection{Matrixschreibweise für komplexe Zahlen
\label{buch:komplex:subsection:matrixschreibweise}}
In der Konstruktion der Körpererweiterung $\mathbb{Q}(\alpha)$
mit einer Nullstelle $\alpha$ des Minimalpolynoms $m(X)\in\mathbb{Q}[X]$ 
wurde nie die Frage gestellt, woher $\alpha$ eigentlich kommt.
Für $m=X^2-2$ kann sich $\alpha$ als die Quadratwurzel
$\!\sqrt{2}\in\mathbb{R}$ vorstellen, aber für die Körpererweiterung
von $\mathbb{F}_5$ um eine Nullstelle von $X^2-2\in\mathbb{F}_5[X]$
gibt es keinen solchen Überkörper, der alle denkbaren Wurzeln
enthält.
Und bei der Erweiterung $\mathbb{R}\subset \mathbb{R}(i)=\mathbb{C}$
wird $\mathbb{R}$ die imaginäre Einheit $i$ hinzugefügt, für die es
ebenfalls keine Realisierung in $\mathbb{R}$ gibt.
Schlimmer noch: die imaginäre Einheit hat Eigenschaften, die den 
Eigenschaften aller bisher bekannten Zahlen widerspricht, nämlich
die Eigenschaft, dass das Quadrat einer Zahl immer $\ge 0$ ist.

So schreibt auch Leonhard Euler in seiner \emph{Vollständigen
Anleitung zur Algebra}
\index{Euler, Leonhard}%
\cite[S.~61,  Abschnitt~144--145]{buch:euler-algebra}:
\begin{quote}
Daher bedeuten alle diese Ausdrücke $\!\sqrt{-1}$, $\!\sqrt{-2}$,
$\!\sqrt{-3}$, $\!\sqrt{-4}$, etc. solche unmögliche oder
Imaginärzahlen, weil dadurch Quadratwurzeln von Negativzahlen
angezeigt werden.

Von diesen behauptet man also mit allem Rechte, dass sie weder
grösser oder kleiner sind als nichts; und auch nicht einmal nichts
selbst, als aus welchem Grund sie folglich für unmöglich gehalten
werden müssen.

Gleichwohl aber werden sie unserm Verstande dargestellet, und finden
in unserer Einbildung statt; daher sie auch bloss eingebildete
Zahlen genennet werden.
Ungeachtet aber diese Zahle, als z.~B. $\!\sqrt{-4}$, ihrer Natur
nach ganz und gar unmöglich sind, so haben wir davon doch einen
hinlänglichen Begriff, indem wir wissen, dass dadurch eine solche
Zahl angedeutet werde, welche mit sich selbst multiplicirt zum
Producte -4 hervorbringe; und dieser Begriff ist zureichend um diese
Zahlen in der Rechnung gehörig zu behandeln.
\end{quote}
Euler war sich also des Widerspruchs durchaus bewusst und versucht
auch gar nicht, eine ``physikalische'' Rechtfertigung für die
imaginären Zahlen zu konstruieren.
Vielmehr weisst er darauf hin, dass es genügt, dass wir mit diesen
Zahlen ohne Widerspruch algebraisch arbeiten können, so wie wir
das bei der Konstruktion der algebraischen Körpererweiterung getan
haben.

Trotzdem muss man die Frage stellen, wie Euler sicher sein konnte,
dass die Algebra der komplexen Zahlen wirklich widerspruchsfrei ist.
Seine umfassende Erfahrung und Intuition in der Algebra dieser Zahlen
wird ihn wohl zu dieser Annahme verleitet haben, ohne dass er einen
weiteren Gedanken darauf verwendet, wie man diese Widerspruchsfreiheit
beweisen könnte.
Zugegebenermassen ist dies für die komplexen Zahlen auch nicht eine
grosse Schwierigkeit, da sich mit der Darstellung von $\mathbb{C}$
als komplexe Zahlenebene eine Konstruktion als Paare von reellen
Zahlen mit Rechenoperationen, für die man die Axiome eines Körpers
verifizieren kann, aufdrängt.

Etwas weniger klar ist die Situation jedoch bei algebraischen
Erweiterungen von $\mathbb{Q}$ und endlichen Körpern.
Zwar ist für Nullstellen von quadratischen Polynomen
immer noch die Vorstellung einer Ebene angemessen, für Körpererweiterungen
um Nullstellen von Polynomen beliebigen Grades funktioniert diese
Intuition jedoch nicht mehr.
Gesucht ist daher eine allgemeine Konstruktion eines Modells für 
algebraische Körpererweiterungen $K(\alpha)$ für Nullstellen beliebiger
Minimalpolynome $m(X)\in K[X]$.
Im Folgenden soll gezeigt werden, wie ein solches Modell innerhalb
der Matrizenalgebra konstruiert werden kann.

\subsubsection{Algebraische Erweiterungen als Matrizen}
Wir beginnen mit einem Beispiel.
Sei $A$ die Matrix
\begin{equation}
A 
=
\begin{pmatrix}
0&2\\
1&0
\end{pmatrix}
\in
M_{2\times 2}(\mathbb{Q}).
\label{buch:komplex:zahlen:eqn:Awurzel2}
\end{equation}
Das Quadrat dieser Matrix ist
\[
A^2
=
\begin{pmatrix}
0&2\\
1&0
\end{pmatrix}
\begin{pmatrix}
0&2\\
1&0
\end{pmatrix}
=
\begin{pmatrix}
2&0\\
0&1
\end{pmatrix}
=
2I
\qquad\text{oder}\qquad
A^2-2I = 0,
\]
wobei $I$ die Einheitsmatrix ist.
Mit $A$ ist also ein algebraisches Objekt gefunden, welches die
Polynomgleichung $x^2-2=0$ erfüllt, sofern man den konstanten
Term des Polynoms auf der linken Seite als $2A^0=2I$ liest und 
die 0 auf der rechten Seite als die Nullmatrix.

Die inverse Matrix von $A$ ist
\[
A^{-1}
=
\begin{pmatrix}
0       & 1 \\
\frac12 & 0
\end{pmatrix}
=
\frac12 A,
\]
was sich mit der früher gefundenen Formel für die Inverse eines
algebraischen Elementes deckt.
Die negative Matrix $-A$ ist eine weitere Nullstelle der Gleichung
$x^2-2=0$.

Die Konstruktion ist natürlich auf den Körper $\mathbb{Q}$ oder
die Konstante $2$ beschränkt.
Die Matrix 
\[
A
=
\begin{pmatrix}
0 & a \\
1 & 0
\end{pmatrix}
\qquad\text{mit Inverser}\qquad
A^{-1}
=
\begin{pmatrix}
0      & 1 \\
a^{-1} & 0
\end{pmatrix}
=
a^{-1}A
\]
ist eine Nullstelle der Gleichung $x^2-a=0$ für ein beliebiges 
invertierbares Element $a\in K^*$ des Koeffizientenkörpers.
$-A$ ist eine weitere solche Lösung.

Die Beispiele zeigen, dass die Matrizenalgebra offenbar flexibel
genug ist, um für beliebige Quadratwurzeln eine Matrixdarstellung
zu finden.
Die Darstellung ist allerdings nicht eindeutig, wie der Versuch
der Lösung der quadratische Gleichung
\[
x^2 + px + q = 0
\]
zeigt.
Nach der bekannten Lösungsformel erwarten wir die Lösung in der
Form
\[
x_{\pm}
=
-\frac{p}2 \pm \!\sqrt{\rlap{$\phantom{\bigg|}$}\smash{\biggl(\frac{p}{2}\biggr)^2 - q}}.
\]
Übersetzt in die Matrixform müsste 
\[
X_{\pm}
=
-\frac{p}2I
\pm
\begin{pmatrix}
0 & \sqrt{(p/2)^2-q} \\
1 & 0
\end{pmatrix}
\]
eine Lösung der Matrixgleichung 
\begin{equation}
X^2 + pX + q = 0
\label{buch:komplex:zahlen:eqn:pqgleichung}
\end{equation}
sein.
Wir berechnen dazu das Produkt
\begin{align*}
X_{\pm}^2
&=
\begin{pmatrix}
(p/2)^2 - q & 0 \\
      0     & (p/2)^2 - q
\end{pmatrix}
=
((p/2)^2-q) I
\end{align*}
und setzen in die Gleichung
\[
X_\pm^2 + pX + q
=
((p/2)^2-q)I
+
pX_\pm
+
qI
=
\begin{pmatrix}
p^2/4 & q + p^3/4 -pq \\
p + q  & p^2/4
\end{pmatrix}
\]
ein.
Die Gleichung ist also nicht erfüllt.
Der Grund ist, dass es in der Menge $M_{2\times 2}(K)$ unendlich
viele Quadratwurzeln gibt.
Ist $T\in \operatorname{GL}_2(K)$ eine beliebige invertierbare 
$2\times 2$-Matrix, dann gilt für $A'=TAT^{-1}$
\[
A^{\prime 2}
=
(TAT^{-1})(TAT^{-1})
=
TA(\underbrace{T^{-1}T}_{\displaystyle=I})AT^{-1}
=
TA^2T^{-1}.
\]
Ist $A$ eine Quadratwurzel von $a$, dann ist
\[
A^{\prime 2}
=
TA^2T^{-1}
=
T(aI)T^{-1}
=
aI,
\]
$A'$ ist also auch eine Quadratwurzel von $a$.
Die Matrizenalgebra ist also viel grösser als zur Darstellung von
Quadratwurzeln nötig ist, zu gross um eine Quadratwurzel eindeutig 
festzulegen.

%
% Eine Matrix als Nullstelle des Minimalpolynoms
%
\subsubsection{Eine Matrix als Nullstelle des Minimalpolynoms}
Es lässt sich aber trotzdem eine Matrix angeben, die als Lösung
der Gleichung \eqref{buch:komplex:zahlen:eqn:pqgleichung} angesehen
werden kann, nämlich
\[
X
=
\begin{pmatrix}
-p & -q \\
 1 &  0
\end{pmatrix}.
\]
Tatsächlich ist
\begin{align*}
X^2 + pX + qI
&=
\begin{pmatrix}
-p & -q \\
 1 &  0
\end{pmatrix}^2
+
p
\begin{pmatrix}
-p & -q \\
 1 &  0
\end{pmatrix}^2
+qI
\\
&=
\begin{pmatrix}
p^2 - q & pq \\
   -p   & -q
\end{pmatrix}
+
\begin{pmatrix}
-p^2 & -pq \\
 p   &  0 
\end{pmatrix}
+
\begin{pmatrix}
q & 0 \\
0 & q
\end{pmatrix}
=
\begin{pmatrix}
0 & 0 \\
0 & 0
\end{pmatrix}.
\end{align*}
Die Matrix $X$ ist also eine Lösung der quadratischen Gleichung
\eqref{buch:komplex:zahlen:eqn:pqgleichung}.

Die Konstruktion des letzten Beispiels lässt sich verallgemeinern
zum folgenden Satz, der besagt, dass jedes über $K$ algebraische
Element vom Grad $n$ eine Matrixdarstellung in $M_{n\times n}(K)$
hat, die Nullstelle der Polynomgleichung ist, die das algebraische
Element definiert.

\begin{satz}
\label{buch:komplex:zahlen:satz:matrixalgebraisch}
Ist $a(X) = X^n + a_{n-1}X^{n-1}+\dots + a_1X+a_0 \in K[X]$ ein
normiertes Polynom vom Grad $n$ mit Koeffizienten im Körper $K$,
dann ist die Matrix
\[
A
=
\begin{pmatrix}
-a_{n-1} & -a_{n-2} & \dots  &  -a_1   &  -a_0   \\
    1    &     0    & \dots  &    0    &    0    \\
    0    &     1    & \dots  &    0    &    0    \\
 \vdots  &  \vdots  & \ddots &  \vdots &  \vdots \\
    0    &     0    & \dots  &    1    &    0 
\end{pmatrix}
\]
eine Nullstelle des Polynoms, also $a(A)=0$.
\end{satz}

\begin{proof}
Wir berechnen als erstes das charakteristische Polynom von $X$ mithilfe
des laplaceschen Entwicklungssatzes, indem wir nach der ersten Spalte
entwickeln.
Zur Vereinfachung schreiben wir
\[
A_{k}
=
\begin{pmatrix}
    -a_{k-1}     & -a_{k-2} & \dots  &  -a_1   &  -a_0   \mathstrut\\
        1        & -\lambda & \dots  &    0    &    0    \mathstrut\\
        0        &     1    & \dots  &    0    &    0    \mathstrut\\
     \vdots      &  \vdots  & \ddots &  \vdots &  \vdots \\
        0        &     0    & \dots  &    1    & -\lambda\mathstrut
\end{pmatrix}
\in
M_{k\times k}(K).
\]
Für die Determinante von $A_k$ folgt durch Entwicklung nach der ersten
Spalte die Rekursionsformel
\begin{align}
\det A_k
&=
-a_{k-1}
\left|
\begin{matrix}
 -\lambda &     0    & \dots  &    0     &    0    \mathstrut\\
     1    & -\lambda & \dots  &    0     &    0    \mathstrut\\
  \vdots  &  \vdots  & \ddots &  \vdots  &  \vdots \\
     0    &     0    & \dots  & -\lambda &    0    \mathstrut\\
     0    &     0    & \dots  &    1     & -\lambda\mathstrut
\end{matrix}
\right|
-
\left|
\begin{matrix}
 -a_{k-2} & -a_{k-3} & \dots  &  -a_1    &  -a_0   \mathstrut\\
     1    & -\lambda & \dots  &    0     &    0    \mathstrut\\
  \vdots  & \vdots   & \ddots &  \vdots  &  \vdots \\
     0    &    0     & \dots  & -\lambda &    0    \mathstrut\\
     0    &    0     & \dots  &    1     & -\lambda\mathstrut
\end{matrix}
\right|
\notag
\\
&=
-a_{k-1}(-\lambda)^{k-1}
-\det A_{k-2}
\label{buch:komplex:zahlen:eqn:xkrekursion}
\intertext{und die Anfangsbedingung}
\det A_2
&=
\biggl|
\begin{matrix}
-a_1 & -a_0     \mathstrut\\
  1  & -\lambda \mathstrut
\end{matrix}
\biggr|
=
a_1\lambda + a_0.
\label{buch:komplex:zahlen:eqn:xkstart}
\end{align}
Durch wiederholte Anwendung der Rekursionsformel
\eqref{buch:komplex:zahlen:eqn:xkrekursion}
ausgehend vom Anfangswert
\eqref{buch:komplex:zahlen:eqn:xkstart}
\begin{align*}
\det A_3
&=
-a_2(-\lambda)^2 - \det A_2
&&=
-a_2\lambda^2 -a_1\lambda - a_0
\\
\det A_4
&=
-a_3(-\lambda)^3 - \det A_3
&&=
a_3 \lambda^3 + a_2\lambda^2 + a_1\lambda + a_0
\\
\det A_5
&=
-a_4(-\lambda)^4 - \det A_4
&&=
-a_4\lambda^4
-a_3\lambda^3 - a_2\lambda^2 -a_1\lambda - a_0
\\
&\phantom{n}\vdots
\\
\det A_k
&=
-a_{k-1}(-\lambda)^{k-1}
-\det A_{k-1}
&&=
(-1)^k\bigl(
a_{k-1}\lambda^{k-1} + \dots + a_1\lambda + a_0
\bigr)
\end{align*}
findet man Ausdrücke für die Determinante $\det A_k$.

Damit kann man jetzt auch das charakteristische Polynom bestimmen.
Dieses ist
\begin{align}
\chi_A(\lambda)
&=
\det (A-\lambda I)
\notag
\\
&=
(-a_{n-1}-\lambda)
\left|
\begin{matrix}
 -\lambda & \dots  &    0    &    0    \mathstrut\\
     1    & \dots  &    0    &    0    \mathstrut\\
  \vdots  & \ddots &  \vdots &  \vdots \mathstrut\\
     0    & \dots  &    1    & -\lambda\mathstrut
\end{matrix}
\right|
-
\left|
\begin{matrix}
-a_{n-2} & -a_{n-3} & \dots  &  -a_1   &  -a_0   \mathstrut\\
    1    & -\lambda & \dots  &    0    &    0    \mathstrut\\
 \vdots  &  \vdots  & \ddots &  \vdots &  \vdots \mathstrut\\
    0    &     0    & \dots  &    1    & -\lambda\mathstrut
\end{matrix}
\right|
\notag
\\
&=
(-a_{n-1}-\lambda)(-\lambda)^{n-1}
-
\det A_{n-2}
\\
&=
(-1)^{n-2}
\bigl(
\lambda^n
+
a_{n-1}\lambda^{n-1}
\bigr)
-
(-1)^{n-1}
\bigl(
a_{n-2}\lambda^{n-2} + \dots + a_1\lambda + a_0
\bigr)
\notag
\\
&=
(-1)^n
\bigl(
\lambda^n
+
a_{n-1}\lambda^{n-1}
+
a_{n-2}\lambda^{n-2}
+
\dots
+
a_1\lambda
+
a_0
\bigr)
\notag
\\
&=
(-1)^n a(\lambda).
\label{buch:komplex:zahlen:eqn:charmin}
\end{align}
Das charakteristische Polynom ist also bis auf ein Vorzeichen das
ursprüngliche Polynom $a(X)\in K[x]$.

Der Satz von Cayley-Hamilton \cite[Satz 13.23]{buch:linalg} besagt, dass 
\index{Cayley-Hamilton, Satz von}%
\index{Satz!von Cayley-Hamilton}%
jede Matrix Nullstelle ihres charakteristischen Polynoms ist, dass also
$\chi_A(A)=0$ ist.
Es folgt daher aus \eqref{buch:komplex:zahlen:eqn:charmin},
dass auch $m(X)=0$ ist.
Damit ist der Satz vollständig bewiesen.
\end{proof}

\begin{beispiel}
Für das irreduzible Polynom $a(X) = X^2-2$ ist
$a_1=0$ und $a_0=-2$.
Die Formel des Satzes liefert daher für 
\[
A
=
\begin{pmatrix}
-a_1&-a_0\\
  1 &  0
\end{pmatrix}
=
\begin{pmatrix}
0 & 2 \\
1 & 0
\end{pmatrix},
\]
also genau die Matrix, die in
\eqref{buch:komplex:zahlen:eqn:Awurzel2}
als Quadratwurzel von $2$ vorgeschlagen wurde.
\end{beispiel}

%
% Die imaginäre Einheit als Matrix
%
\subsubsection{Die imaginäre Einheit als Matrix}
Die imaginäre Einheit ist Nullstelle des irreduziblen Polynoms
$X^2+1$.
Nach Satz~\ref{buch:komplex:zahlen:satz:matrixalgebraisch} kann
die imaginäre Einheit durch die $2\times 2$-Matrix
\[
J
=
\begin{pmatrix*}[r]
0 & -1 \\
1 &  0
\end{pmatrix*}
\]
dargestellt werden.
Tatsächlich sind die Potenzen von $J$
\begin{align*}
J^2
&=
\begin{pmatrix*}[r]
-1& 0\\
 0&-1
\end{pmatrix*}
=
-I
\\
J^3
&=
J^2J
=
-IJ
=
-J
\\
J^4
&=
J^2J^2
(-I)(-I)
=
I.
\\
J^k
&=
\begin{cases}
(-1)^{k/2}I&\qquad \text{$k$ gerade} \\
(-1)^{(k-1)/2}J&\qquad \text{$k$ ungerade,} \\
\end{cases}
\end{align*}
wie man dies von $i$ erwartet.

Die komplexe Zahl $z=a+bi$ wird also durch die Matrix
\[
Z = aI+bJ
=
\begin{pmatrix*}[r]
a&-b\\
b& a
\end{pmatrix*}
\]
dargestellt.
Wir rechnen nach, dass die inverse Matrix von $Z$ auch die
Matrix ist, die zur komplexen Zahl $z^{-1}$ gehört.
Die inverse Matrix ist
\begin{align*}
Z^{-1}
&=
\frac{1}{\det Z}
\begin{pmatrix*}[r]
 a & b \\
-b & a
\end{pmatrix*}
=
\frac{1}{a^2+b^2}
\begin{pmatrix*}[r]
 a & b \\
-b & a
\end{pmatrix*}
=
\frac{a}{a^2+b^2}I
+
\frac{-b}{a^2+b^2}J.
\end{align*}
Dies ist genau die Matrix, die zur reziproken komplexen Zahl
\[
\frac{1}{z}
=
\frac{a}{a^2+b^2}
+
\frac{-b}{a^2+b^2}i
\] 
gehört.

Man beachte, dass in dieser Matrixdarstellung der komplexen Zahlen
die erste Spalte ein Vektor ist, der die komplexe Zahl als Vektor
in der gaussschen Zahlenebene darstellt.
